% !TEX root = report.tex
\S4.1 is analytical. The rest use 200 seeds per configuration.
Unless stated, Red plays MIP-BR with $\varepsilon_\Rd=0.15$ in nominal weather.

\subsection{Template game and a small example: patrol or assert, and sustainable restraint}
IESDS reduces $\Gamma$ from $14\times14$ to $2\times4$. Blue keeps only the center-heavy split (20/60/20), and Red
keeps the south-heavy and center/south-heavy splits. Of the two pure NE, the \emph{equilibrium} policy plays (Blue center-heavy
assert, Red south-heavy patrol), $(0.39,-9.27)$, over (Blue patrol, Red center/south assert), $(-14.44,2.71)$.
Without the escalation cost, the zero-sum maximin LP has a pure equilibrium where \emph{both} sides concentrate on
the Center and assert. The shared incident cost is what makes at most one side assert.
Fixing Blue center-heavy and Red south-heavy, each side only chooses P or A. The resulting $2\times2$ sub-matrix of $\Gamma$, normalized so (P,P) $=(0,0)$, has
(P,A) $=(-9.61,\,1.20)$, (A,P) $=(0.28,\,-10.45)$ and (A,A) $=(-11.41,\,-11.33)$ (Blue, Red).
Asserting against a patroller pays a little (one-period gain $G_\B=0.28$, $G_\Rd=1.20$), but mutual assertion costs more than
patrolling alone, so $BR_i(\Pt)=\As$ and $BR_i(\As)=\Pt$. There are two pure NE, (A,P) and (P,A), and a mixed NE (Blue asserts w.p.\
$0.58$, Red $0.13$) with mutual assertion w.p.\ $0.08$ and payoffs $v_\B=-1.30$, $v_\Rd=-6.02$. The welfare-maximizing correlated equilibrium (CE LP)
puts \emph{zero} on (A,A). A mediator (e.g.\ a hotline) removes mutual assertion.

In the repeated game, suppose both patrol, and any assertion makes both fall back to the
mixed NE for good. By the one-stage deviation principle, patrolling is optimal iff
$(1-\delta)G_i+\delta v_i\le0$, i.e.\ $\delta\ge G_i/(G_i-v_i)$, which is $0.18$ for Blue and $0.17$ for Red. With
$\delta=0.98$, \textbf{mutual restraint is a subgame perfect equilibrium}, though our one-period-ahead policies
cannot see this.

\subsection{RQ1: control versus war risk, and the risk-budget game}
\noindent\begin{minipage}[t]{0.40\textwidth}\vspace{0pt}\centering\footnotesize\setlength{\tabcolsep}{2pt}
\begin{tabular}{l|ccc}
\toprule
$\B\backslash\Rd$ & cons.\ $0.05$ & mod.\ $0.15$ & aggr.\ $0.40$\\\midrule
cons.\ $0.05$ & $-194,\ 199$ & $-303,\ 315$ & $-308,\ 318$\\
              & war $0.10$ & $0.11$ & $0.12$\\
mod.\ $0.15$  & $18,\ -36$ & $-144,\ 101$ & $-192,\ 140$\\
              & war $0.15$ & $0.20$ & $0.21$\\
aggr.\ $0.40$ & $74,\ -80$ & $-83,\ 16$ & $\mathbf{-151,\ 39}$\\
              & war $0.14$ & $0.25$ & $\mathbf{0.34}$\\
\bottomrule
\end{tabular}
\captionof{table}{The risk-budget game: $(u_\B,u_\Rd)$ and war rate when both sides choose $\varepsilon$
(s.e.\ $\approx 5$--$12$). Bold: the unique NE.}\label{fig:rq1}
\end{minipage}\hfill
\begin{minipage}[t]{0.58\textwidth}\vspace{0pt}
\textbf{Frontier:} Raising $\varepsilon_\B$ from $0.02$ to $0.15$ lifts Blue's control share from $-0.25$ to
$-0.05$ while the war rate rises from $0.07$ to $0.20$ (Red fixed at $0.15$). Going to $0.40$ adds only $0.06$ of control. Most of what a looser budget buys is bought by $\varepsilon\approx0.15$--$0.25$.

\smallskip\textbf{Risk-budget game} (Table~\ref{fig:rq1})\textbf{:} each side picks a budget. Whatever the other
side picks, ``aggressive'' pays each side the most, so both pick it, and
(aggressive, aggressive) is the unique pure NE. It also has the \emph{highest} war rate ($0.34$), and both sides
would do at least as well if both chose moderate ($-144,\,101$ vs.\ $-151,\,39$, war $0.20$). This is a prisoner's
dilemma, so restraint must come from repetition or coordination (\S4.1).
\end{minipage}

\subsection{RQ2: strategic versus heuristic play}
\begin{table}[!ht]\centering\footnotesize\setlength{\tabcolsep}{4pt}\renewcommand{\arraystretch}{0.9}
\begin{tabular}{l|ccc}
\toprule
Blue $\backslash$ Red & rule-based (moderate) & MIP-BR & equilibrium\\\midrule
random & $-283$ / $-0.14$ / $0.15$ & $-319$ / $-0.16$ / $0.33$ & $-164$ / $-0.08$ / $0.06$\\
rule-based (moderate) & $-122$ / $-0.05$ / $0.04$ & $-184$ / $-0.09$ / $0.06$ & $35$ / $+0.06$ / $0.04$\\
greedy & $-70$ / $+0.01$ / $0.12$ & $-137$ / $-0.05$ / $0.14$ & $65$ / $+0.08$ / $0.03$\\
MIP-BR & $-75$ / $+0.01$ / $0.13$ & $-144$ / $-0.05$ / $0.20$ & $53$ / $+0.08$ / $0.08$\\
equilibrium & $-72$ / $-0.01$ / $0.07$ & $\mathbf{-116}$ / $-0.03$ / $0.10$ & $21$ / $+0.05$ / $0.06$\\
maximin & $-170$ / $-0.09$ / $0.04$ & $-221$ / $-0.12$ / $0.08$ & $-35$ / $-0.01$ / $0.00$\\
\bottomrule
\end{tabular}
\caption{Blue payoff $u_\B$ / control share / war rate (s.e.\ of $u_\B$ is 1--7). Nominal weather, design D0. The
MIP-BR column is each policy's worst case.}
\label{tab:policies}
\end{table}
Greedy, MIP-BR and equilibrium beat the random and rule-based benchmarks (Table~\ref{tab:policies}). The
\textbf{equilibrium} policy has the best worst case ($-116$) and, against MIP-BR, half the war rate of MIP-BR itself
($0.10$ vs.\ $0.20$). Greedy and MIP-BR do best against weaker opponents but escalate more. \textbf{Maximin} is the
safest (war rate $\le0.08$) but gives up the most control, guarding against templates Red rarely plays.
Greedy matched the MIP optimum in 67\% of 1{,}524 Blue decisions but, unlike the MIP, broke the risk budget in 14\%.

\subsection{RQ3: robustness to weather and resilient designs}
% !TEX root = updates.tex
\noindent\begin{minipage}[t]{0.47\textwidth}\vspace{0pt}\centering\footnotesize\setlength{\tabcolsep}{2pt}
{\begin{tabular}{l|cccc|cc}
\toprule
Design & calm & nom. & typh. & extr. & exp. & worst\\\midrule
D0 baseline   & $-458$/.34 & $-430$/.26 & $-399$/.21 & $-366$/.16 & $-423$ & $-458$\\
D1 dispersed  & $-474$/.34 & $-421$/.23 & $-394$/.19 & $-368$/.11 & $-422$ & $-474$\\
D2 sensing    & $-479$/.34 & $-447$/.28 & $-424$/.28 & $-386$/.12 & $-443$ & $-479$\\
D3 de-escalation & $\mathbf{-444}$/.10 & $\mathbf{-395}$/.11 & $\mathbf{-382}$/.06 & $\mathbf{-350}$/.04 & $\mathbf{-400}$ & $\mathbf{-444}$\\
\bottomrule
\end{tabular}}
\captionof{table}{Blue designs by weather scenario (MIP-BR vs.\ MIP-BR): $u_\B$ / war rate (s.e.\ $\approx6$--$9$).}\label{fig:rq3}
\smallskip\normalsize\raggedright We score each policy or design $x$ by its expected payoff
$\sum_\omega\pi_\omega J(x,\omega)$ and worst case $\min_\omega J(x,\omega)$.
\end{minipage}\hfill
\begin{minipage}[t]{0.51\textwidth}\vspace{0pt}
\textbf{Weather:} as in Workshop 1, bad weather helps Blue. Every policy except MIP-BR with $\varepsilon=0.40$ does
worst in \emph{calm} seas, since rough seas halve small-hull strength, which is about half of Red's strength but less than a quarter of
Blue's, and storms
pause operations and incidents (MIP-BR war rate $0.34$ calm, $0.16$ extreme). The top of the ranking is the same in all
four scenarios: MIP-BR with $\varepsilon=0.40$, then $\varepsilon=0.15$, then greedy. Only the three weakest policies
(rule-based, equilibrium, maximin) change places. The best policy on average ($\varepsilon=0.40$, expected $3$) is also
best in its worst case ($-88$, extreme seas).
\end{minipage}

\textbf{Designs:} de-escalation measures (D3) are again best in expected and worst-case terms. They cut the war rate
from $0.26$ to $0.11$ (nominal) \wone{$0.20\to0.05$} and raise Blue's payoff by $35$. Red's payoff rises too
($385\to417$), since fewer wars help both sides. We still charge them nothing, so part of this win is by assumption.
Dispersed basing (D1) leaves Blue's payoff where it was ($-421$ against $-430$) and no longer cuts Red's
\wone{$101\to16$}, now $385\to371$. Redundant sensing (D2) is worse than the baseline ($-447$), paying its cost for no
measurable benefit.
\textbf{Limitations:} parameters are assumptions, policies look one period ahead, and each game ship stands for three
real ships. Payoff differences below about $20$ are not resolved with 200 games (s.e.\ $6$--$13$).

